\section{Discussion and Limitations}
\label{sec:discussion}

Excluding the Rat snapshot preserves the main average gains over the static and endpoint-only forecasts, as well as the endpoint's advantage over the alternative summaries on the joint bases. These observations reduce concern that the input-construction result depends on the 78 unrecovered Rat preprocessing decisions. They do not establish that those decisions were appropriate or harmless within Rat. We retain that dataset's original results for traceability and treat them as evidence conditional on the archived snapshot; a claim of robustness to Rat preprocessing would require a separate comparison starting from explicitly prepared data and refitting the bases.

The input ablations strengthen the case for combining the endpoint with compressed history and forecast coefficients: the complete construction improves average error over endpoint-only regression, including without blending. They also delimit that claim. Retaining residual DCT inputs slightly worsens MSE on some ETTh2--DLinear conditions, and the complete regression uses more storage than its reduced-input controls. The evidence supports average complementarity in the evaluated setting rather than a universally necessary input set.

The setting and temporal-selection analyses similarly support robustness of the observed average benefit, not an independent confirmation. Larger $K$ values improve mean accuracy at additional storage cost, and prefix selection frequently chooses $K=8$. The reference $K=4$ is a compression choice. Although selection uses only preceding losses within each replay, the series and candidate construction have already informed development. Neither retrospective splitting nor a small gap to the finite-pool oracle removes that broader dependence. The archived selection history is partial, so these diagnostics cannot quantify the entire research process's selection bias.

The endpoint ablation suggests that selecting information to retain is a useful design choice beyond selecting the dimension of a compressed representation. The first, middle, and endpoint variants have the same regression dimensions and reconstruct corrections in the same DCT subspace. Their differences therefore arise from the supplied residual value and the parameters estimated from it, rather than from an increase in model size. A compressed block summary and a separately retained observation can play complementary roles. The former describes variation across the block, while the latter preserves a value that the retained coefficients need not determine.

This interpretation does not require the endpoint to contain information absent from the full residual block used by a learned adapter. Its role is to remain explicitly available after compression. The comparison with the first and middle residuals supports choosing the most recent value within the tested construction. It does not establish which physical process makes that value predictive. Likewise, the larger gains over learned adapters involve differences in estimation and architecture as well as representation. The equal-size ablation provides the more direct evidence for the endpoint choice.

The adverse interval in Fig.~\ref{fig:worst_interval} has a different implication: a correction can improve average accuracy while becoming temporarily harmful. The regression discounts older blocks, and the controller estimates correction strength from a finite window of completed blocks. Both therefore depend on past performance. If the relationship between past residuals and future errors changes, a previously useful correction may remain active before subsequent observations reduce its influence. This is a plausible consequence of the update mechanism, but the observed trace alone does not establish it as the cause of the selected failure.

The distinction matters when interpreting the bounded blending coefficient. Restricting it to $[0,1]$ limits the amplitude applied to the raw correction. It does not bound the resulting error relative to the base forecast on an unseen target block. Average error reduction and local excess error consequently describe different properties of the method. For applications in which short adverse intervals matter, the duration and magnitude of such intervals should inform the assessment of correction, alongside full-period accuracy. An explicit control on local harm would require an additional design and evaluation objective beyond the retrospective squared-error fit used here.

Finally, the storage comparison measures retained numerical arrays rather than hardware performance. Peak runtime memory also depends on input buffers, temporary arrays, and the solver implementation. Latency and energy consumption depend on the execution platform and computational workload. We have not measured these quantities on a target device, so the reduction in retained arrays cannot be translated directly into a measured reduction in RAM, latency, or energy. Evaluating that relationship requires a separate device implementation and measurement study.
